\paragraph{Learning from the first image the writer actually drew.}
At a known insertion boundary, changing the complete first glyph can change every later character through visual feedback. We retained the previous writer's continuation and stop parameters and fitted two first-drawing corrections to the same 241,256 training sites. One maximized reference image likelihood. The second also learned to prefer the recorded glyph to the frozen writer's own competing first image.

On 24,048 held-out insertions, likelihood training increased exact first images from 4,920 to 4,970 and exact stopped phrases from 2,143 to 2,177. Preference training produced 5,090 exact first images and 2,280 exact stopped phrases (Fig.~\ref{fig:first-emission}). It corrected 678 erroneous first images, lost 508 previously correct first images and changed another 7,214 into a different wrong image. Against 19,128 previous first-image errors, the reference signature now outranked the previous drawing in 1,089 cases. The model recovered the reference bitmap in 631 of those cases. In 336 cases a newly emitted competing signature still outranked the reference.

The preference writer recovered 148 of 13,041 multi-glyph insertions and one of 7,504 insertions absent from the expanded training inventory. Whole-phrase counts gave 300 exact multi-glyph and 3,554 total insertions. On the separately reserved chapters the preference model recovered 33 of 837 phrases, the frozen writer 26 and counts 59. The chapter-bootstrap interval for the preference gain in complete answers over the frozen writer was 0.23 to 1.59 percentage points. Against the count control it was $-4.87$ to $-1.25$ points. A blinded contextual reading accepted six of 48 preference answers, four likelihood answers, four frozen answers, eleven count answers and all 48 references. A single study-involved model reader made these judgments. In 797 natural source-exchange pairs, preference training favored the authentic assignment in 459 pairs. Actual generation recovered 173 of 1,594 phrases with the authentic classical source and 158 after exchange. These experiments establish a limited improvement in visually drawn Chinese insertions with a supplied modern prefix.

Removing both supplied inputs exposed a larger failure. The editor now chose its own insertion boundaries and carried its own rendered history through the entire sentence. Across all 84 reserved sentences, neither writer recovered a complete reference. The preference model stopped on 22 and reached its 32-glyph phrase limit on 62. Its glyph edit distance was 5,532 over 5,994 reference glyphs, compared with 5,576 for the frozen writer and 4,578 from copying the classical source. Exact reference matching penalizes other acceptable translations, but the repeated and incomplete outputs also limit practical use. Figure~\ref{fig:first-emission}d reproduces the complete output of a short sentence.

Of the 62 preference outputs ending at the phrase limit, 52 ended in an exact bitmap sequence repeated at least three times consecutively. None of their recorded target sentences had such a suffix. We also tested whether an offline text teacher could resolve the underlying source ambiguity before considering distillation. In a calibration sample of 32 pairs of different multi-glyph insertions drawn from matched local contexts, Qwen3 4B selected the authentic phrase at 24 of 64 sites under both option orders, falling to five with the classical sources exchanged. The original A/B parser left 71 of 256 responses undecided. Resolving unambiguous Chinese option names and requesting a constrained response for the remaining 40 replies raised the counts to 35 and eight. The source supplied information, but neither pass reached the prespecified 40 correct sites. The visual preference writer ranked the authentic candidate above the alternative at 33 sites with authentic sources and 31 after exchange. Candidate recognition requires much less than independently completing a translation.

Using the complete translations as candidates made this same assignment task much easier. Across all 6,720 calibration pairs, cosine similarity between bags of identical glyph images correctly matched 6,714 complete sentence pairs to their translations. Adjacent glyph pairs alone solved 6,664, with 51 ties. Only four distinct sentence pairs defeated both overlap rules. Success on this full-sentence assignment roster would therefore give little evidence for reading beyond shared glyphs.
